% ch-linearsolver.tex -- numerical method: the two linear-solver backends and their shared contract
\chapter{Linear solvers}
\label{ch:linearsolver}

Every Newton iteration and every semi-implicit solve ends in a sparse linear
system (\cref{sec:time:newton}). Yuzen ships two implementations of that one
job, and which of them you get is an input-file decision. This chapter says
what the two are, how they compare on measured cases, and what every backend
is required to do regardless of which one you pick. The keys themselves are
tabulated in \cref{sec:ref-numerics:linearsolver}.

\section{Two backends behind one interface}
\label{sec:linsolve:backends}

The two implementations sit behind the single interface
\code{ILinearSolverBackend} and are interchangeable. The calling code knows
only \code{LinearSolver} and never learns which backend it has. The choice is
made by \key{Numerics.HeightEquation.LinearSolver.Backend} and
\key{Numerics.ConcEquation.LinearSolver.Backend}, separately --- the thickness
equation and the concentration equation may use different backends in the same
run.

\begin{yzntable}{The three values of \key{Backend}\label{tab:backends}}{@{}l L{0.30\textwidth} L{0.38\textwidth}@{}}
\textsf{\bfseries Value} & \textsf{\bfseries Implementation} &
\textsf{\bfseries Character} \\ \midrule
\val{"auto"} (default) & amgcl for a 2-D film on a build with AMGCL, otherwise eigen &
  the choice to make when you have no reason to prefer another \\
\val{"eigen"} & Eigen \code{BiCGSTAB} with block-Jacobi \code{IncompleteLUT} &
  depends on Eigen alone; parallel, and \textbf{bit-identical whatever the thread count} \\
\val{"amgcl"} & AMGCL \code{bicgstab} with smoothed-aggregation AMG and SPAI0 &
  fastest in 2-D on most sizes; OpenMP-parallel including the preconditioner \\
\end{yzntable}

\begin{note}[Why there is no in-house BiCGStab]
  An in-house BiCGStab with ILU(0) was measured during development and left
  out of the release: it was far slower than the alternatives --- on
  \file{leveling\_1d\_reg} it needed 3676 linear iterations where
  \val{"eigen"} needs 21 (\cref{sec:verification:backends}) --- and it had
  bugs, while the one property it offered, bit-identity under a change of
  thread count, is a property of \val{"eigen"} as well. Any other value of
  \key{Backend} is \textbf{refused rather than quietly given something
  else}: the parser reports that \key{Backend} \code{must be one of
  [`auto',`eigen',`amgcl']} and the run aborts with exit code 1
  (\cref{sec:ref-numerics:linearsolver}).
\end{note}

Asking for \val{"amgcl"} on a build without AMGCL is not an error: the solver
emits a warning and falls back to \val{"eigen"}. The line in
\file{Warnings.txt} reads \code{LinearSolver: Backend="amgcl" was requested,
but this build has no AMGCL. Falling back to "eigen".} Grep for
\code{has no AMGCL}. The run then produces valid
results that are not the ones you asked for, so check which backend was
actually used before comparing against another machine's output. You do not
have to infer it: the \code{\#\#\#\#\#\# Solver Info \#\#\#\#\#\#} section of
\file{JobInfo.txt} names the resolved backend of each equation
(\cref{sec:output:jobinfo}), so it reads \val{"eigen"} after a fall-back and
names the real choice behind \val{"auto"}.

\section{What each backend does}
\label{sec:linsolve:what}

\subsection{\val{"eigen"}: block-Jacobi ILUT}
\label{sec:linsolve:eigen}

With \key{PreconditionType} \val{"ILU"}, the rows are split into contiguous
blocks, one per 1024 rows and at most 128, and each diagonal block is
factorized by Eigen's \code{IncompleteLUT} (fill factor 10, drop tolerance
$10^{-10}$); the preconditioner applies each block's factors to its own rows,
and the coupling between the blocks is left to the BiCGStab iteration. The
blocks are factorized and applied in parallel. \textbf{The blocking follows
from the size of the system alone, never from the thread count}, so the
result is bit-identical whatever the thread count, as before
(\cref{ch:reproducibility}). A system of fewer than 2048 rows is one block:
the ILUT of the whole matrix. Against one ILUT of the whole matrix, the blocks
take 4 to 9 times the iterations on the 2-D films measured, and are 1.3 to 3
times faster on one thread all the same, because the factorization of the
whole matrix costs far more than that of its blocks; at the best thread
count, 3.7 to 21 times faster (\cref{sec:parallel:profile}). The analysis of each block's
pattern is done once and reused while the pattern stays the same.

\subsection{\val{"amgcl"}: smoothed-aggregation AMG}
\label{sec:linsolve:amgcl}

With \val{"ILU"}, AMGCL's BiCGStab is preconditioned by smoothed-aggregation
AMG with SPAI0 smoothing, built anew at every solve. On a system of more
than 3000 unknowns the hierarchy is coarsened down to 500 (AMGCL's own default
is 3000), because the coarsest level is solved directly and sequentially: on a
$128^2$ film the deeper hierarchy made the solve 2.9 times faster. A system of
3000 unknowns or fewer is left to AMGCL's default, which solves it directly in
one iteration; coarsened to 500, the fourth-order 1-D film of 2000 cells took
80 BiCGStab iterations a solve and ran 6.5 times slower on one thread.

\subsection{The concentration equation}
\label{sec:linsolve:conc}

The concentration systems are dominated by their diagonal (the film over the
step), so with \key{PreconditionType} \val{"ILU"} a light preconditioner is
tried first: Jacobi with \val{"eigen"}, SPAI0 with \val{"amgcl"}, at most 300
iterations. \textbf{If that try does not
converge, the solve starts again from the same initial guess with the
preconditioner of the input} (ILUT or AMG). A system on which the light one
does not work --- a ConcZ film with dry columns is one --- then costs a few
failed tries, not one per solve: after a failure the next 16 tries are
skipped, doubling after each further failure up to 1024, and back to 16 after
a try that converged. The same holds for the block-Jacobi ILUT of the film
height and of the concentration: a try that does not converge in 300
iterations falls back to the ILUT of the whole matrix. Which way a solve goes
depends on its convergence only, not on the thread count. The iterations of a
dropped try are counted in the linear iterations of \file{JobInfo.txt}.

\subsection{The tolerance of each Newton pass}
\label{sec:linsolve:ew}

The film height is solved by Newton's method (\cref{sec:time:newton}), one
linear solve a pass. A pass far from the solution does not need its
correction to 13 digits, so \textbf{the tolerance of each pass follows the
progress of the iteration} (Eisenstat and Walker's choice 2):
\[
  \eta_k = 0.9\left(\frac{\lVert F_k\rVert}{\lVert F_{k-1}\rVert}\right)^{2},
\]
$F$ the right-hand side of the pass (the residual of the step at the current
guess), safeguarded by $0.9\,\eta_{k-1}^2$ while that is above $0.1$, and held
between \key{Epsilon} and $10^{-4}$; the first pass of a step is solved to
$10^{-4}$. \textbf{\key{Epsilon} of the film height's linear solver is
therefore the tightest tolerance a pass is solved to}, reached as the
iteration converges, not the tolerance of every pass. The Newton convergence
tests themselves (\cref{sec:convergence:tests}) are unchanged, and the
answer agrees with that of solving every pass to \key{Epsilon} to within
the Newton tolerance: on the cases measured, the final film differed by
$10^{-15}$ relative and the linear iterations fell by 24 to 38\,\%.

These four choices can each be turned off with \env{YUZEN\_OPT}
(\cref{tab:envvars}), to compare a run with the algorithm it replaced.

\section{Which one to choose}
\label{sec:linsolve:choosing}

\Cref{tab:linsolve-timing} gives measured wall-clock times. They were taken on
one node with two Intel Xeon CPU Max 9480 (56 cores each, 8 NUMA nodes), with
GCC 11 and LLVM's \file{libomp} (the default build, \cref{sec:parallel:runtime}),
from a \code{Release} build configured with \code{-DYUZEN\_WITH\_AMGCL=ON},
the threads bound to cores 0 to $T-1$ (\env{OMP\_PROC\_BIND=close}) and the
shortest of three runs shown --- the same machine as the case run times of
\cref{tab:cases-runtimes}. \textbf{The absolute values depend on the machine,
so what to read is the ratio between the columns, not the seconds
themselves.} \file{tools/case\_runtimes.py --backends-only} measures the table
(\cref{sec:env:tools}).

\begin{yzntable}{Measured wall-clock times, one Xeon Max 9480 node, up to the cores of one socket (\cref{sec:parallel:cap}).\label{tab:linsolve-timing}}{@{}L{0.24\textwidth} l d d@{}}
\textsf{\bfseries Case} & \textsf{\bfseries Threads} &
\textsf{\bfseries eigen} & \textsf{\bfseries amgcl} \\ \midrule
1-D \file{leveling\_1d\_reg}, $N_x = 300$ & 1 & \textbf{0.58 s} & 0.64 s \\
                                         & 2 & \textbf{0.46 s} & 0.51 s \\
                                         & 4 & \textbf{0.38 s} & 0.44 s \\
                                         & 8 & \textbf{0.35 s} & 0.42 s \\
                                         & 16 & \textbf{0.36 s} & 0.42 s \\
                                         & 28 & \textbf{0.36 s} & 0.43 s \\
                                         & 56 & \textbf{0.37 s} & 0.44 s \\
2-D \file{leveling\_2d\_reg}, $60\times60$ & 1 & 19.8 s & \textbf{12.0 s} \\
                                         & 2 & 10.6 s & \textbf{7.5 s} \\
                                         & 4 & 6.5 s & \textbf{4.8 s} \\
                                         & 8 & 5.8 s & \textbf{3.6 s} \\
                                         & 16 & 5.5 s & \textbf{3.1 s} \\
                                         & 28 & 5.4 s & \textbf{2.9 s} \\
                                         & 56 & 5.3 s & \textbf{2.8 s} \\
\end{yzntable}

\textbf{In 1-D eigen is \SI{10}{\percent} faster on one thread and stays
ahead; in 2-D amgcl is 1.6 times as fast on one thread and 1.9 times on 56.}
These cases are small --- 300 and 3\,600 cells. In 1-D both gain a factor of
about 1.6 by eight threads and nothing after; with the runtime's hot teams
(\cref{sec:parallel:runtime}) more threads do not cost time either, 56 being
within \SI{6}{\percent} of 8. In 2-D eigen gains a factor of 3.8 by 56 threads
and amgcl 4.3, most of it by eight. \Cref{sec:parallel:scaling} shows the
larger films.

Where the time of the 2-D case goes (\file{tools/\bk profile\_threads.py},
\cref{sec:parallel:profile}): with eigen, building the block-Jacobi ILUT
preconditioner takes \SI{55}{\percent} of the time-step loop on one thread and
on eight, and the BiCGStab iterations \SI{13}{\percent}; the factorization is
redone at every linear solve, and its blocks stop dividing the work at a few
thousand rows. The profile of amgcl is spread out: on one thread the AMG
set-up takes \SI{9}{\percent}, the iterations \SI{27}{\percent} and the
assembly of the film-height matrices \SI{27}{\percent}.

The in-house backend was timed during development on a two-core Xeon at
\SI{2.1}{\giga\hertz} on the same two cases, and two things in those measurements are worth
remembering. In 1-D it took 9.3 seconds at one thread and 11.5 at two, against
2.2 and 1.9 for eigen on the same development form of the case --- \textbf{four
to six times slower, and \emph{slower} again when given a second thread}. That was not a measurement error: its ILU
had to run sequentially (\cref{sec:repro:other}), so the only parts that
parallelized were the matrix--vector products and the reductions, and at this
problem size the synchronization cost more than those parts saved. In 2-D, at
5.2 and 4.0 seconds, it did beat the 12.9 and 11.4 seconds eigen took on
that development form of the case, but it did not beat amgcl, which is the
backend to reach for in that regime anyway.

\textbf{If exact agreement between runs matters, use \val{"eigen"}.}
AMGCL's internal reductions are not order-fixed, so changing
the thread count can change the low-order bits. Measured on \file{leveling\_1d\_reg} at
one thread against four, the differences are $6.089\times10^{-11}$ in the
physical quantities and $9.615\times10^{-2}$ in the error measures (a column
that is rounding in that case); at two threads, and in 2-D
(\file{allphys\_2d\_reg}) at two and four, the bits happened not to move,
which nothing guarantees. For \val{"eigen"} the same comparison gives exactly
$0$ at any thread count (\cref{sec:verification:backends}). \Cref{ch:reproducibility} explains what
stands behind that property, and why it is worth having.

\begin{warn}[Changing the backend changes how the rounding falls]
  Results from two different backends do not agree bit for bit. They agree only
  to within the convergence tolerance of the linear solver. When you make a
  regression comparison, fix the backend on both sides of it --- and check what
  each run really used in the \code{Solver Info} section of its
  \file{JobInfo.txt} (\cref{sec:output:jobinfo}) rather than trusting the two
  input files to have resolved \val{"auto"} the same way.
\end{warn}

\section{The shared contract}
\label{sec:linsolve:contract}

Whichever backend is in use, the following hold.

\begin{itemize}
  \item \textbf{The solution vector passed to \code{Solve()} is also the
    initial guess.} With Eigen it is handed to \code{solveWithGuess}. The
    return value is \val{0} for converged, \val{1} for the iteration cap
    reached, and \val{2} for divergence or breakdown.
  \item \textbf{A failed solve must have its solution vector thrown away.} If a
    NaN is left in the vector and passed to the next call, the solve fails with
    a NaN again no matter how far the time step is cut, and the run grinds down
    to \key{MinDt} and aborts for a reason that has nothing to do with the time
    step. \code{C\_NewtonKantorovich::Solve} discards it; anything new that
    calls a backend directly must do the same.
  \item \textbf{Diagonal scaling does not change the solution.}
    \key{RowScaling} divides each row by its own diagonal, which leaves the
    solution of the system untouched. It is applied to a mirror of the values
    and to a copy of the right-hand side; the matrix itself is not modified.
  \item \textbf{Preconditioner names map differently onto each
    implementation.} Under \val{"eigen"}, \val{"ILU"} is block-Jacobi ILUT
    (\cref{sec:linsolve:eigen}); under \val{"amgcl"} it becomes
    smoothed-aggregation AMG with SPAI0, and \val{"Jacobi"} becomes damped
    Jacobi. The name you write selects a preconditioner of comparable
    strength, not the same algorithm.
  \item \textbf{\env{YUZEN\_ILUT\_FILL} overrides the fill factor} of Eigen's
    \code{IncompleteLUT}; the default is 10. Lowering it makes the
    preconditioner cheaper, but drop too much and the factorization breaks down
    and the solution comes back as NaN --- a value of 1 breaks down even on a
    1-D problem.
\end{itemize}
